\documentclass[12pt,usepdftitle=false,aspectratio=1610]{beamer}

\usetheme{Boadilla}
%\usetheme{Madrid}
%\usecolortheme{seahorse}

\setbeamertemplate{navigation symbols}{}

%\usefonttheme[onlymath]{serif}
\usefonttheme[]{serif}
%\useinnertheme{circles}
\setbeamertemplate{caption}[numbered]

\usepackage[T2A]{fontenc}			
\usepackage[utf8]{inputenc}			
\usepackage[english,russian]{babel}
\usepackage{textcomp}
\usepackage{gensymb}
\usepackage{mathtext,cmap,multirow} 				
\usepackage{graphicx,wrapfig,subfig,mathtools,gensymb}
\usepackage[font=small,labelfont=bf]{caption}
\usepackage{minted}

\graphicspath{{./figs/01}}

\title[Лекция 1]{Цифровое изображение. Цветовые модели}

\author{Александр Танченко}
\institute{}
\date{2025}

\begin{document}

%------------------------------------------------------------
\begin{frame}
    \titlepage
\end{frame}

%------------------------------------------------------------
\begin{frame}
\frametitle{Черно-белое (binary) изображение}
\begin{figure}
    \centering
    \includegraphics[height=0.3\textheight]{coordinates.png}
    \includegraphics[height=0.3\textheight]{black_whight.png}
\end{figure}
$$
  f(m,n)\in\left\{0,1\right\}
$$
\begin{itemize}
    \item Черно-белое изображение - это одноканальное однобитное изображение.
    \item $0$ - черный цвет, $1$ - белый цвет
    \item хорошо подходит для хранения масок
\end{itemize}
\end{frame}

%------------------------------------------------------------
\begin{frame}
\frametitle{Полутоновое (grayscale) изображение}
\begin{figure}
    \centering
    \includegraphics[height=0.3\textheight]{coordinates.png}
    \includegraphics[height=0.3\textheight]{gray_image.png}
\end{figure}
$$
  f(m,n)\in\left\{0,1,\ldots,255\right\}
$$
\begin{itemize}
    \item Полутоновое изображение - это одноканальное восьмибитное изображение.
    \item При изменении значения пикселя $f(m,n)$ от 0 до 255, 
    цвет меняется в градациях серого от черного к белому.
\end{itemize}
\end{frame}

%------------------------------------------------------------
\begin{frame}
\frametitle{Цветное RGB изображение}
\begin{figure}
    \centering
    \includegraphics[height=0.3\textheight]{RGB_image.png}
    \includegraphics[height=0.3\textheight]{lena.png}
\end{figure}
$$
  f_k(m,n)\in\left\{0,1,\ldots,255\right\}\,\quad
  k=1,2,3
$$
\begin{itemize}
    \item Цветное RGB изображение - это трехканальное восьмибитное изображение.
    \item Каналы $k=1,2,3$ соответствуют красному (R), зеленому (G) и синему (B) цвету.
\end{itemize}
\end{frame}

%------------------------------------------------------------
\begin{frame}
\frametitle{Примеры цифровых изображений}
\begin{figure}
    \centering
    \includegraphics[height=0.3\textheight]{galaxy.jpg}
    \includegraphics[height=0.3\textheight]{thermography.jpg}
\end{figure}
\begin{figure}
    \centering
    \includegraphics[height=0.3\textheight]{mri.jpg}
    \includegraphics[height=0.3\textheight]{electron_microscopy.jpg}
\end{figure}
\end{frame}

%------------------------------------------------------------
\begin{frame}
\frametitle{Математическое описание цифровых изображений}
\begin{figure}
    \centering
    \includegraphics[height=0.3\textheight]{digital_image.jpg}
\end{figure}

Цифровое изображение задается набором $K$ матриц (каналов) $M\times N$ 
и числом бит $B$
$$
    \mathbf{F}=\left\{f_{k}(m,n)\right\},\quad
    m=0,1,\ldots,M-1,\quad
    n=0,1,\ldots,N-1,\quad
    k=1,2,\ldots,K
$$
$$
    f_{k}(m,n)\in\left\{0,1,\ldots,2^B-1\right\}
$$
\end{frame}


%------------------------------------------------------------
\begin{frame}
\frametitle{Формирование цифрового изображения}
\begin{figure}
    \centering
    \includegraphics[height=0.4\textheight]{acquisition.pdf}
\end{figure}
\begin{itemize}
    \item при формировании цифрового изображения непрерывный сигнал преобразуется
    в дискретный
    \item затем дискретный сигнал сжимается и сохраняется
    \item на каждом шаге обработки теряется информация, 
    что приводит к возникновению артефактов
\end{itemize}
\end{frame}

%------------------------------------------------------------
\begin{frame}
\frametitle{Сэмплирование и квантизации непрерывного сигнала}
\begin{figure}
    \centering
    \includegraphics[height=0.55\textheight]{sampling.pdf}
\end{figure}
\begin{itemize}
    \item при сэмплировании значения непрерывного сигнала вычисляются в узлах сетки
    \item при квантизации диапазон значений сигнала  разбивается на конечное число уровней и значения округляются до ближайших к ним уровням
\end{itemize}
\end{frame}

%------------------------------------------------------------
\begin{frame}
\frametitle{Артефакты сэмплирования}
\begin{figure}
    \centering
    \includegraphics[height=0.75\textheight]{sampling_artefacts.pdf}
    \caption*{Артефакты сэмплирования при 930 dpi, 300 dpi, 150 dpi и 75 dpi.}
\end{figure}
\end{frame}

%------------------------------------------------------------
\begin{frame}
\frametitle{Артефакты квантования}
\begin{figure}
    \centering
    \includegraphics[height=0.7\textheight]{quantization_artefacts.pdf}
    \caption*{Шумы камеры и сатурация - артефакты квантизации (слишком большой шаг квантизации).}
\end{figure}
\end{frame}

%------------------------------------------------------------
\begin{frame}
\frametitle{Артефакты квантования}
\begin{figure}
    \centering
    \includegraphics[height=0.5\textheight]{jpeg_artifacts.png}
\end{figure}
График демонстрирует артефакты при JPEG сжатии изображения и
зависимость метрики качества \textbf{PSNR (dB)}
$$
  \mbox{PSNR}(A,B)=10\cdot\log_{10}\frac{255^2\cdot M\cdot N}{||A - B||_F^2}
$$
от коэффициента сжатия.
\end{frame}

%------------------------------------------------------------
\begin{frame}
\frametitle{Видимый свет}
\begin{figure}
    \centering
    \includegraphics[height=0.4\textheight]{light.png}
\end{figure}
\begin{itemize}
    \item \textbf{видимый свет} -- поток электромагнитных волн разных длин волн
    \item \textbf{монохромный свет} - свет состоящий из волн только одной длины волны
    \item человек воспринимает длины волн $360\,\mbox{нм} < \lambda < 830\,\mbox{нм}$
    \item человек отличает примерно $1000$ оттенков цветов и $20$ оттенков серого
\end{itemize}
\end{frame}

%------------------------------------------------------------
\begin{frame}
\frametitle{Спектральная плотность излучения}
\begin{figure}
    \centering
    \includegraphics[height=0.4\textheight]{D65.png}
    \caption*{Cпектральная плотность для дневного света (источник $D65$)}
\end{figure}
\begin{itemize}
    \item свет состоит из множества волн с разными длинами волн $\lambda$
    \item \textbf{спектральная плотность излучения} $I(\lambda)$ --
  плотность потока излучения, приходящуюся на длину волны $\lambda$
\end{itemize}
\end{frame}

%------------------------------------------------------------
\begin{frame}
\frametitle{Элементы зрительного восприятия}
\begin{figure}
    \centering
    \includegraphics[height=0.25\textheight]{eye.pdf}
    \includegraphics[height=0.25\textheight]{absorption.pdf}
\end{figure}
Сетчатка человеческого глаза содержит два вида рецепторов:
\begin{itemize}
    \item \textbf{палочки} -- отвечают за сумеречное зрение
    \item \textbf{колбочки} -- обеспечивают цветовое зрение
    \item cуществует три типа колбочек, которые в основном чувствительны к синему, зеленому и красному свету
\end{itemize}
\end{frame}

%------------------------------------------------------------
\begin{frame}
\frametitle{Эксперимент Райта и Гилда}
\begin{figure}
    \centering
    \includegraphics[height=0.45\textheight]{color_experiment.jpg}
\end{figure}
Эксперимент Дэвида Райта и Джона Гилда в 1920-х годах:
\begin{itemize}
    \item многие цвета нельзя разложить в сумму трех базовых цветов
    \item если к цвету, который не получается разложить, добавить определенный цвет из
    базовых, то получившийся цвет раскладывается по трём базовым цветам
\end{itemize}
\end{frame}

%------------------------------------------------------------
\begin{frame}
\frametitle{Цветовая модель XYZ}
\begin{figure}
    \centering
    \includegraphics[height=0.3\textheight]{XYZ.png}
\end{figure}
В 1931 году разработана цветовая модель \textbf{XYZ}:
\begin{itemize}
    \item зафиксированы три положительные весовые функции
    $x(\lambda)$, $y(\lambda)$, $z(\lambda)$
    \item координаты света со  спектральной плотностью $I(\lambda)$ равны
    $$
        X = \int\limits_{360}^{830} {x}(\lambda)I(\lambda)\,d\lambda\,,\quad
        Y = \int\limits_{360}^{830} {y}(\lambda)I(\lambda)\,d\lambda\,,\quad
        Z = \int\limits_{360}^{830} {z}(\lambda)I(\lambda)\,d\lambda
    $$
\end{itemize}
\end{frame}

%------------------------------------------------------------
\begin{frame}
\frametitle{Цветовая модель XYZ}
\begin{figure}
    \centering
    \includegraphics[height=0.5\textheight]{CIEXYZ.png}
\end{figure}
Из формул видно, что отображение $I(\lambda) \mapsto (X\,,Y\,,Z)$ линейно:
\begin{itemize}
    \item при сложении световых потоков координаты цветов тоже складываются
    \item при увеличении потока в $\alpha$ раз его координаты тоже увеличиваются
    в $\alpha$ раз
\end{itemize}
\end{frame}

%------------------------------------------------------------
\begin{frame}
\frametitle{Цветовая модель xyY}
Удобно перейти к так называемым \textbf{хроматическим координатам} $x$ и $y$ при помощи следующих уравнений:
$$
    x = \frac{X}{X + Y + Z}\,,\quad
    y = \frac{Y}{X + Y + Z}
$$

В результате мы получаем цветовую модель \textbf{xyY}.

Ниже приводятся формулы для обратных преобразований из \textbf{xyY} в \textbf{XYZ}:
$$
    X = \frac{Y}{y}x\,,\quad
    Z = \frac{Y}{y}(1 - x - y)
$$
\end{frame}

%------------------------------------------------------------
\begin{frame}
\frametitle{Хроматическая диаграмма}
\begin{figure}
    \centering
    \includegraphics[height=0.5\textheight]{Yxy.png}
\end{figure}
\begin{itemize}
    \item \textbf{xроматическая диаграмма} -- множества хроматических координат $(x,y)$
    всех различаемых глазом человека цветов
    \item дуга на диаграмме соответствует монохроматическим цветам $360\,\mbox{нм} < \lambda  < 830\,\mbox{нм}$
    \item все цвета внутри любого треугольника можно представить в виде взвешенной суммы трех цветов в вершинах треугольника
\end{itemize}
\end{frame}

%------------------------------------------------------------
\begin{frame}
\frametitle{Цветовая модель sRGB}
\begin{figure}
    \centering
    \includegraphics[height=0.45\textheight]{sRGB.png}
\end{figure}
$$
    \begin{bmatrix}
    R \\ G \\ B
    \end{bmatrix} =
    \begin{bmatrix}
        3.2405 & -1.5372 & -0.4985\\
        -0.9693 & 1.8760& 0.0416\\
        0.0556 & -0.2040 & 1.0573
    \end{bmatrix}
    \cdot
    \begin{bmatrix}
    X \\ Y \\ Z
    \end{bmatrix}
$$
$$
    R =
    \begin{cases}
        1.055 \cdot R^{1/2.4} - 0.055\,, & R > 0.0031\\
        12.92 \cdot R\,,                 & R \leq 0.0031
    \end{cases}
    \,,\quad
    R =
    \begin{cases}
        0\,,  & R < 0\\
        1 \,, & R > 1
    \end{cases}
$$

Для $G$, $B$ формулы аналогичные.
\end{frame}

%------------------------------------------------------------
\begin{frame}
\frametitle{Свойства RGB}
\begin{figure}
    \centering
    \includegraphics[height=0.45\textheight]{RGB.pdf}
    \includegraphics[height=0.35\textheight]{RGB_cube.pdf}
\end{figure}
\begin{itemize}
    \item не все цвета могут быть представлены в \textbf{RGB} модели
    \item модель \textbf{RGB}, как и \textbf{XYZ}, является аддитивной (при сложении световых потоков соответствующие координаты цветов тоже складываются)
\end{itemize}
\end{frame}

%------------------------------------------------------------
\begin{frame}
\frametitle{Цветовое пространство HSV}
\begin{figure}
    \centering
    \includegraphics[height=0.35\textheight]{HSV.png}
\end{figure}
\textbf{HSV} моделирует, как художники воспринимают цвет
\begin{itemize}
    \item $0^\circ \leqslant H < 360^\circ$  \textbf{тон (hue)}, чистый оттенок цвета
    \item $0\leqslant S \leqslant 1$  \textbf{насыщенность (saturation)},
    насколько цвет близок к чистому цвету
    \item $0< V \leqslant 1$  \textbf{яркость (value)}
\end{itemize}
\end{frame}

%------------------------------------------------------------
\begin{frame}
\frametitle{Цветовое пространство HSV}
\begin{figure}
    \centering
    \includegraphics[height=0.35\textheight]{HSV.png}
\end{figure}
$$
    V = \max(R,G,B)\,,\quad
    S=1-\frac{\min(R,G,B)}{\max(R,G,B)}
$$
$$
    h =
    \begin{cases}
        60(G-B)/\delta\,, & R=V\\
        60(B-R)/\delta+120\,, & G=V\\
        60(R-G)/\delta+240\,, & B=V
    \end{cases}\,,\quad
    \delta = V(1-S)
$$
$$
    H = h \mod 360^\circ
$$
\end{frame}

%------------------------------------------------------------
\begin{frame}
\frametitle{Цветовая модель YCbCr}
\begin{figure}
    \centering
    \includegraphics[height=0.3\textheight]{YCbCr.png}
\end{figure}
\begin{align*}
    Y  &= 0.299\cdot R + 0.587\cdot G + 0.114\cdot B \\
    Cb &= 0.564 \cdot (B - Y) + 0.5 \\
    Cr &= 0.713 \cdot (R - Y) + 0.5
\end{align*}
\begin{itemize}
    \item $Y$ -- канал яркости (Grayscale)
    \item $Cb$, $Cr$ -- цветовые каналы
\end{itemize}
\end{frame}

%------------------------------------------------------------
\begin{frame}
\frametitle{Цветовая аугментация изображений. Изменение тона}
\begin{figure}
    \centering
    \includegraphics[height=0.4\textheight]{hue_adjust.png}
\end{figure}
\begin{itemize}
    \item сконвертировать изображение в цветовое пространство \textbf{HSV}
    \item изменить плоскость тона \textbf{H} (hue) по формуле
    $$
      \mathbf{H} = (\mathbf{H} + \alpha) \mod 360\,,\quad
      0 < \alpha < 360
    $$
    \item сконвертировать изображение обратно в цветовое пространство \textbf{RGB}
\end{itemize}
\end{frame}

%------------------------------------------------------------
\begin{frame}
\frametitle{Изменение насыщенности}
\begin{figure}
    \centering
    \includegraphics[height=0.4\textheight]{saturation_adjust.png}
\end{figure}
\begin{itemize}
    \item найти яркость изображения $\mathbf{Y}$
    \item преобразовать цветовые плоскости \textbf{R}, \textbf{G} и \textbf{B} по формулам
    \begin{align}
      \mathbf{R} &= \beta \cdot \mathbf{R} + (1-\beta)\cdot\mathbf{Y}\\
      \mathbf{G} &= \beta \cdot \mathbf{G} + (1-\beta)\cdot\mathbf{Y}\\
      \mathbf{B} &= \beta \cdot \mathbf{B} + (1-\beta)\cdot\mathbf{Y}
  \end{align}
  где $0 < \beta < 1$
\end{itemize}
\end{frame}

%------------------------------------------------------------
\begin{frame}
\frametitle{Изменение контрастности}
\begin{figure}
    \centering
    \includegraphics[height=0.4\textheight]{contrast_adjust.png}
\end{figure}
\begin{itemize}
    \item найти яркость изображения $\mathbf{Y}$ и вычислить среднюю яркость $\bar{\mathbf{Y}}$
    \item преобразовать цветовые плоскости \textbf{R}, \textbf{G} и \textbf{B} по формулам
    \begin{align*}
      \mathbf{R} &= \gamma \cdot \mathbf{R} + (1-\gamma)\cdot\bar{\mathbf{Y}}\\
      \mathbf{G} &= \gamma \cdot \mathbf{G} + (1-\gamma)\cdot\bar{\mathbf{Y}} \\
      \mathbf{B} &= \gamma \cdot \mathbf{B} + (1-\gamma)\cdot\bar{\mathbf{Y}}
  \end{align*}
  где $0 < \gamma < 1$
\end{itemize}
\end{frame}

%------------------------------------------------------------
\begin{frame}
\frametitle{Изменение яркости}
\begin{figure}
    \centering
    \includegraphics[height=0.4\textheight]{brightness_adjust.png}
\end{figure}
\begin{itemize}
    \item цветовые плоскости \textbf{R}, \textbf{G} и \textbf{B} умножить на коэффициент $0 < \delta < 1$
    \begin{align*}
        \mathbf{R} &= \delta \cdot \mathbf{R} \\
        \mathbf{G} &= \delta \cdot \mathbf{G} \\
        \mathbf{B} &= \delta \cdot \mathbf{B}
    \end{align*}
\end{itemize}
\end{frame}

%------------------------------------------------------------
\begin{frame}
\frametitle{ColorJitter аугментация}
\begin{itemize}
    \item случайным образом выбрать коэффициенты $\alpha$, $\beta$, $\gamma$, $\delta$ из диапазона
    \begin{align*}
        0 < \alpha_1 &< \alpha < \alpha_2 < 360 \\
        0 < \beta_1 &< \beta < \beta_2 < 1 \\
        0 < \gamma_1 &< \gamma < \gamma_2 < 1 \\
        0 < \delta_1 &< \delta < \delta_2 < 1
    \end{align*}
    \item к изображению последовательно применить цветовые преобразования
    тона, насыщенности, контрастности и яркости
    с параметрами $\alpha$, $\beta$, $\gamma$ и $\delta$
\end{itemize}
\end{frame}

\end{document}

